\section{Implementation and Reproducibility}
\label{app:reproducibility}

\paragraph{Environment.}
Each task ran in a Docker environment connected to self-hosted GitLab,
RocketChat, OwnCloud, and Plane services. We used \texttt{pi-agent} v0.75.3 as
the fixed executor harness and exposed bash as its only action
tool. All service interaction, including API calls, was therefore
issued through shell commands. Every attempt was launched with
\texttt{--no-session}, so the harness wrote no session state to disk and
nothing it might otherwise retain could carry across tasks or rounds.
Self-hosted services were confirmed running before each execution batch.

\paragraph{Environment reset between attempts.}

\paragraph{Models and retrieval.}
Claude Opus 4.6, GPT-5.4, and GLM-5 are the executor-consolidator model
variants. Claude Sonnet 4.6 performs evaluation, and Cohere Embed v4
provides 1024-dimensional semantic-search embeddings. The LLM wiki was available through link navigation from
\texttt{index.md} and semantic search over page contents. All conditions used
the same executor loop and a retrieval stage before planning. Consolidated
representations used retrieval-only agents with representation-specific
adapters, while \cond{No Memory} omitted retrieval and \cond{No Consolidation} supplied raw
rollout evidence. Table~\ref{tab:retrieval-adapters} gives the exact
protocols.

\paragraph{Evaluator protocol for checkpoint scoring.}

\paragraph{Compute, time, and human effort.}
Experiments were orchestrated on an Amazon EC2 \texttt{g5.2xlarge} host with
8 vCPUs and 30 GiB of memory. Model inference was served through
provider APIs,
while the host ran the agent harness, benchmark services, and Docker
environments. One complete 127-task evaluation batch occupied approximately 12
wall-clock hours. Producing \cond{Checkpoint-binary + Human} feedback for one rollout
batch required approximately six hours from a qualified annotator.

\begin{table}[t]
\caption{\textbf{API cost of one complete 127-task evaluation batch.}
Execution is the task-execution phase, retrieval is the memory-query phase, and
\cond{Free Form} total is their sum. Consolidation costs are reported separately in
the experiment tables. Evaluator, embedding, host-compute, and
human-feedback costs are outside these values.}
\label{tab:execution-cost}
\centering
\small
\resizebox{\linewidth}{!}{%
\begin{tabular}{lrrrr}
\toprule
Model & No Memory execution & Free Form execution & Retrieval &
Free Form total \\
\midrule
Claude Opus 4.6 & \$106.6 & \$88.9 & \$60.4 & \$149.3 \\
GPT-5.4 & \$108.5 & \$87.1 & \$60.6 & \$147.8 \\
GLM-5 & \$56.3 & \$47.2 & \$18.6 & \$65.9 \\
\bottomrule
\end{tabular}}
\end{table}

\paragraph{Round procedure.}
Within each round, all benchmark tasks executed against a frozen memory store and
produced one rollout batch. Feedback records were generated from checkpoint
evaluation and, for the rich condition, human failure diagnosis. The
consolidator then updated a copy of the previous wiki. The semantic index was
rebuilt before the next execute phase. R0 used no store and supplied the rollout
batch from which R1 memory was created.

\paragraph{Validation.}
Metric denominators were normalized to the benchmark totals for every
condition. Shared conditions reuse one canonical result wherever they appear
across experiments. All main and appendix tables, and the values behind
every figure, are transcribed from the same source-of-truth result pages.

\paragraph{Experiment code.}
